% Generated by generate_latex.py
\documentclass[12pt]{article}
\usepackage[utf8]{inputenc}
\usepackage[T2A]{fontenc}
\usepackage[russian]{babel}
\usepackage{microtype}
\usepackage{geometry}
\geometry{a4paper, left=25mm, right=25mm, top=25mm, bottom=25mm}
\usepackage{enumitem}
\usepackage{tcolorbox}
\tcbset{colback=white, colframe=black!30, boxrule=0.6pt, arc=4pt, left=4pt, right=4pt, top=4pt, bottom=4pt}
\usepackage[hidelinks]{hyperref}
\usepackage{titlesec}
\titleformat{\section}{\Large\bfseries}{}{0pt}{}
\begin{document}
\begin{center}
  {\LARGE Сборник вопросов}\\[6pt]
  {\large Сгенерировано автоматически}
\end{center}
\vspace{6mm}
\section*{Введение}
\begin{enumerate}[leftmargin=*]
\end{enumerate}
\vspace{2mm}
\section*{Перечислите способы создания мобильных приложений}
\begin{enumerate}[leftmargin=*]
\end{enumerate}
\vspace{2mm}
\section*{С каким фреймворком используют язык Dart при создании мобильных приложений}
\begin{enumerate}[leftmargin=*]
\end{enumerate}
\vspace{2mm}
\section*{С каким фреймворком используют язык C# при создании мобильных приложений}
\begin{enumerate}[leftmargin=*]
\end{enumerate}
\vspace{2mm}
\section*{На каком языке создают мобильные приложения для iPhone в наиболее стандартном случае}
\begin{enumerate}[leftmargin=*]
\end{enumerate}
\vspace{2mm}
\section*{Опишите особенности функциональной парадигмы программирования}
\begin{enumerate}[leftmargin=*]
\end{enumerate}
\vspace{2mm}
\section*{Перечислите основные парадигмы программирования}
\begin{enumerate}[leftmargin=*]
\end{enumerate}
\vspace{2mm}
\section*{Чем отличается нативная разработка от кроссплатформенной}
\begin{enumerate}[leftmargin=*]
\end{enumerate}
\vspace{2mm}
\section*{Что такое SDK}
\begin{enumerate}[leftmargin=*]
\end{enumerate}
\vspace{2mm}
\section*{Синтаксис}
\begin{enumerate}[leftmargin=*]
\end{enumerate}
\vspace{2mm}
\section*{Опишите синтаксис описания переменных в языке Kotlin}
\begin{enumerate}[leftmargin=*]
\end{enumerate}
\vspace{2mm}
\section*{Опишите синтаксис и семантику оператора for языка Kotlin}
\begin{enumerate}[leftmargin=*]
\end{enumerate}
\vspace{2mm}
\section*{Опишите синтаксис и семантику оператора when языка Kotlin}
\begin{enumerate}[leftmargin=*]
\end{enumerate}
\vspace{2mm}
\section*{Опишите синтаксис описания функции на языке Kotlin}
\begin{enumerate}[leftmargin=*]
\end{enumerate}
\vspace{2mm}
\section*{Опишите синтаксис лямбда-функции на языке Kotlin}
\begin{enumerate}[leftmargin=*]
\end{enumerate}
\vspace{2mm}
\section*{Опишите механизм задания значений параметров функций по умолчанию и правила их переопределения при вызове}
\begin{enumerate}[leftmargin=*]
\end{enumerate}
\vspace{2mm}
\section*{Опишите назначение типа Unit и отличие от void в других языках}
\begin{enumerate}[leftmargin=*]
\end{enumerate}
\vspace{2mm}
\section*{Как объявить и вызвать функцию с переменным числом аргументов}
\begin{enumerate}[leftmargin=*]
\end{enumerate}
\vspace{2mm}
\section*{Объясните, как реализована типобезопасность в сигнатурах функций Kotlin}
\begin{enumerate}[leftmargin=*]
\end{enumerate}
\vspace{2mm}
\section*{Простые вопросы}
\begin{enumerate}[leftmargin=*]
\end{enumerate}
\vspace{2mm}
\section*{Как язык Kotlin поддерживает null-безопасность}
\begin{enumerate}[leftmargin=*]
\end{enumerate}
\vspace{2mm}
\section*{Чем отличается оператор if языка Kotlin от одноименных операторов языков C и Python}
\begin{enumerate}[leftmargin=*]
\end{enumerate}
\vspace{2mm}
\section*{Чем отличается оператор when языка Kotlin от операторов switch других C-подобных языков}
\begin{enumerate}[leftmargin=*]
\end{enumerate}
\vspace{2mm}
\section*{Чем отличается список и массив}
\begin{enumerate}[leftmargin=*]
\end{enumerate}
\vspace{2mm}
\section*{Перечислите основные коллекции языка Kotlin}
\begin{enumerate}[leftmargin=*]
\end{enumerate}
\vspace{2mm}
\section*{Опишите особенности Sequence}
\begin{enumerate}[leftmargin=*]
\end{enumerate}
\vspace{2mm}
\section*{Чем отличается List от Sequence}
\begin{enumerate}[leftmargin=*]
\end{enumerate}
\vspace{2mm}
\section*{Чем отличается List от Set}
\begin{enumerate}[leftmargin=*]
\end{enumerate}
\vspace{2mm}
\section*{Чем отличается Set от Map}
\begin{enumerate}[leftmargin=*]
\end{enumerate}
\vspace{2mm}
\section*{В чем отличие Sequence от обычных коллекций в плане вычислений и производительности?}
\begin{enumerate}[leftmargin=*]
\end{enumerate}
\vspace{2mm}
\section*{Как работает ленивая обработка коллекций в Kotlin?}
\begin{enumerate}[leftmargin=*]
\end{enumerate}
\vspace{2mm}
\section*{Что произойдет, если вызвать toList() у Sequence?}
\begin{enumerate}[leftmargin=*]
\end{enumerate}
\vspace{2mm}
\section*{Написать функцию}
\begin{enumerate}[leftmargin=*]
\end{enumerate}
\vspace{2mm}
\section*{Напишите функцию, которая возвращает функцию, возвращающую последовательно числа 1, 2, 3 и так далее}
\begin{enumerate}[leftmargin=*]
\end{enumerate}
\vspace{2mm}
\section*{Напишите функцию, которая по двум данным функциям Int->Int возвращает их сумму}
\begin{enumerate}[leftmargin=*]
\end{enumerate}
\vspace{2mm}
\section*{Напишите функцию, которая возвращает функцию, возвращающую последовательно элементы массива, а затем null}
\begin{enumerate}[leftmargin=*]
\end{enumerate}
\vspace{2mm}
\section*{Напишите функцию, которая по двум данным функциям Int->Int возвращает их произведение}
\begin{enumerate}[leftmargin=*]
\end{enumerate}
\vspace{2mm}
\section*{Напишите функцию, принимающую список функций Int->Int и возвращающую их композицию}
\begin{enumerate}[leftmargin=*]
\end{enumerate}
\vspace{2mm}
\section*{Напишите функцию, которая возвращает функцию, создающую независимые счетчики с собственным внутренним состоянием}
\begin{enumerate}[leftmargin=*]
\end{enumerate}
\vspace{2mm}
\section*{Напишите функцию, которая принимает Int->Int и возвращает Int->Int, ограничивающую диапазон возвращаемых значений}
\begin{enumerate}[leftmargin=*]
\end{enumerate}
\vspace{2mm}
\section*{Напишите функцию, которая по функции возвращает мемоизированный вариант той же функции}
\begin{enumerate}[leftmargin=*]
\end{enumerate}
\vspace{2mm}
\section*{Напишите функцию, принимающую список функций (Int)->Int и применяющую их последовательно к начальному значению}
\begin{enumerate}[leftmargin=*]
\end{enumerate}
\vspace{2mm}
\section*{Напишите функцию, возвращающую функцию, вычисляющую скользящее среднее последовательности чисел}
\begin{enumerate}[leftmargin=*]
\end{enumerate}
\vspace{2mm}
\section*{ООП}
\begin{enumerate}[leftmargin=*]
\end{enumerate}
\vspace{2mm}
\section*{Опишите особенности реализации объектно-ориентированной парадигмы языком Kotlin}
\begin{enumerate}[leftmargin=*]
\end{enumerate}
\vspace{2mm}
\section*{Как создаются конструкторы в языке Kotlin}
\begin{enumerate}[leftmargin=*]
\end{enumerate}
\vspace{2mm}
\section*{Как описывать свойства классов языка Kotlin}
\begin{enumerate}[leftmargin=*]
\end{enumerate}
\vspace{2mm}
\section*{Какой аналог static членов класса в языке Kotlin}
\begin{enumerate}[leftmargin=*]
\end{enumerate}
\vspace{2mm}
\section*{Опишите синтаксис и семантику делегирования классов}
\begin{enumerate}[leftmargin=*]
\end{enumerate}
\vspace{2mm}
\section*{Опишите синтаксис и семантику делегирования свойств}
\begin{enumerate}[leftmargin=*]
\end{enumerate}
\vspace{2mm}
\section*{Назначение и особенности data class}
\begin{enumerate}[leftmargin=*]
\end{enumerate}
\vspace{2mm}
\section*{Назначение и особенности sealed class}
\begin{enumerate}[leftmargin=*]
\end{enumerate}
\vspace{2mm}
\section*{Особенности enum class языка Kotlin}
\begin{enumerate}[leftmargin=*]
\end{enumerate}
\vspace{2mm}
\section*{Чем отличаются обычные и extension-методы}
\begin{enumerate}[leftmargin=*]
\end{enumerate}
\vspace{2mm}
\section*{Что такое fun interface}
\begin{enumerate}[leftmargin=*]
\end{enumerate}
\vspace{2mm}
\section*{Опишите причину введения в язык отдельных data class}
\begin{enumerate}[leftmargin=*]
\end{enumerate}
\vspace{2mm}
\section*{Что такое object и companion object}
\begin{enumerate}[leftmargin=*]
\end{enumerate}
\vspace{2mm}
\section*{Зачем используется ключевое слово abstract и в каких случаях оно необходимо}
\begin{enumerate}[leftmargin=*]
\end{enumerate}
\vspace{2mm}
\section*{Что такое интерфейсы и как они реализуются}
\begin{enumerate}[leftmargin=*]
\end{enumerate}
\vspace{2mm}
\section*{Как работает наследование}
\begin{enumerate}[leftmargin=*]
\end{enumerate}
\vspace{2mm}
\section*{Что такое переопределение методов}
\begin{enumerate}[leftmargin=*]
\end{enumerate}
\vspace{2mm}
\section*{Что такое вложенные и внутренние классы}
\begin{enumerate}[leftmargin=*]
\end{enumerate}
\vspace{2mm}
\section*{Чем отличаются sealed и enum классы}
\begin{enumerate}[leftmargin=*]
\end{enumerate}
\vspace{2mm}
\section*{Объясните, зачем используется ключевое слово open и как оно влияет на иерархию наследования}
\begin{enumerate}[leftmargin=*]
\end{enumerate}
\vspace{2mm}
\section*{Опишите работу by lazy}
\begin{enumerate}[leftmargin=*]
\end{enumerate}
\vspace{2mm}
\section*{Опишите работу by Delegates.observable}
\begin{enumerate}[leftmargin=*]
\end{enumerate}
\vspace{2mm}
\section*{Что такое lateinit var и почему его не стоит использовать?}
\begin{enumerate}[leftmargin=*]
\end{enumerate}
\vspace{2mm}
\section*{В чём отличие const val от обычного val?}
\begin{enumerate}[leftmargin=*]
\end{enumerate}
\vspace{2mm}
\section*{Функции высшего порядка}
\begin{enumerate}[leftmargin=*]
\end{enumerate}
\vspace{2mm}
\section*{Опишите функцию map}
\begin{enumerate}[leftmargin=*]
\end{enumerate}
\vspace{2mm}
\section*{Опишите функцию zip}
\begin{enumerate}[leftmargin=*]
\end{enumerate}
\vspace{2mm}
\section*{Опишите функцию unzip}
\begin{enumerate}[leftmargin=*]
\end{enumerate}
\vspace{2mm}
\section*{Опишите функцию fold}
\begin{enumerate}[leftmargin=*]
\end{enumerate}
\vspace{2mm}
\section*{Опишите функцию reduce}
\begin{enumerate}[leftmargin=*]
\end{enumerate}
\vspace{2mm}
\section*{Опишите функцию flatten}
\begin{enumerate}[leftmargin=*]
\end{enumerate}
\vspace{2mm}
\section*{Опишите функцию partition}
\begin{enumerate}[leftmargin=*]
\end{enumerate}
\vspace{2mm}
\section*{Опишите функцию groupBy}
\begin{enumerate}[leftmargin=*]
\end{enumerate}
\vspace{2mm}
\section*{Опишите функции take и drop}
\begin{enumerate}[leftmargin=*]
\end{enumerate}
\vspace{2mm}
\section*{Опишите функцию filter}
\begin{enumerate}[leftmargin=*]
\end{enumerate}
\vspace{2mm}
\section*{Опишите функции any и all}
\begin{enumerate}[leftmargin=*]
\end{enumerate}
\vspace{2mm}
\section*{Что делает операция to?}
\begin{enumerate}[leftmargin=*]
\end{enumerate}
\vspace{2mm}
\section*{Что делает функция flatMap, и чем она отличается от map?}
\begin{enumerate}[leftmargin=*]
\end{enumerate}
\vspace{2mm}
\section*{Чем отличаются dropWhile и filter при работе с коллекциями?}
\begin{enumerate}[leftmargin=*]
\end{enumerate}
\vspace{2mm}
\section*{Корутины}
\begin{enumerate}[leftmargin=*]
\end{enumerate}
\vspace{2mm}
\section*{Что такое корутина}
\begin{enumerate}[leftmargin=*]
\end{enumerate}
\vspace{2mm}
\section*{Чем отличаются suspend функции от обычных}
\begin{enumerate}[leftmargin=*]
\end{enumerate}
\vspace{2mm}
\section*{Зачем нужны диспетчеры для корутин, какие бывают диспетчеры}
\begin{enumerate}[leftmargin=*]
\end{enumerate}
\vspace{2mm}
\section*{Как запустить корутину, чтобы из нее можно было получить результат}
\begin{enumerate}[leftmargin=*]
\end{enumerate}
\vspace{2mm}
\section*{Зачем нужны CoroutineScope}
\begin{enumerate}[leftmargin=*]
\end{enumerate}
\vspace{2mm}
\section*{Зачем нужны Flow}
\begin{enumerate}[leftmargin=*]
\end{enumerate}
\vspace{2mm}
\section*{Чем отличаются StateFlow и SharedFlow}
\begin{enumerate}[leftmargin=*]
\end{enumerate}
\vspace{2mm}
\section*{Как отменить корутину и что при этом происходит внутри механизма исполнения}
\begin{enumerate}[leftmargin=*]
\end{enumerate}
\vspace{2mm}
\section*{Что делает функция async и чем она отличается от launch}
\begin{enumerate}[leftmargin=*]
\end{enumerate}
\vspace{2mm}
\section*{Объясните, как обеспечивается упорядоченное выполнение и синхронизация корутин}
\begin{enumerate}[leftmargin=*]
\end{enumerate}
\vspace{2mm}
\section*{Чем отличаются launch, async и runBlocking?}
\begin{enumerate}[leftmargin=*]
\end{enumerate}
\vspace{2mm}
\section*{Что делает функция withContext и зачем она нужна?}
\begin{enumerate}[leftmargin=*]
\end{enumerate}
\vspace{2mm}
\section*{Что произойдёт, если не обработать исключение внутри корутины?}
\begin{enumerate}[leftmargin=*]
\end{enumerate}
\vspace{2mm}
\section*{Чем отличается GlobalScope от пользовательского CoroutineScope?}
\begin{enumerate}[leftmargin=*]
\end{enumerate}
\vspace{2mm}
\section*{Вспомогательные функции и конструкции}
\begin{enumerate}[leftmargin=*]
\end{enumerate}
\vspace{2mm}
\section*{Опишите оператор ?.}
\begin{enumerate}[leftmargin=*]
\end{enumerate}
\vspace{2mm}
\section*{Опишите оператор !!, почему он не рекомендуется к использованию?}
\begin{enumerate}[leftmargin=*]
\end{enumerate}
\vspace{2mm}
\section*{Опишите оператор ?:}
\begin{enumerate}[leftmargin=*]
\end{enumerate}
\vspace{2mm}
\section*{Опишите функцию let}
\begin{enumerate}[leftmargin=*]
\end{enumerate}
\vspace{2mm}
\section*{Опишите функцию run}
\begin{enumerate}[leftmargin=*]
\end{enumerate}
\vspace{2mm}
\section*{Опишите функцию apply}
\begin{enumerate}[leftmargin=*]
\end{enumerate}
\vspace{2mm}
\section*{Опишите функцию also}
\begin{enumerate}[leftmargin=*]
\end{enumerate}
\vspace{2mm}
\section*{Опишите функцию with}
\begin{enumerate}[leftmargin=*]
\end{enumerate}
\vspace{2mm}
\section*{Чем отличаются apply, also, run, let, with}
\begin{enumerate}[leftmargin=*]
\end{enumerate}
\vspace{2mm}
\section*{Как использовать elvis-оператор для задания значения по умолчанию}
\begin{enumerate}[leftmargin=*]
\end{enumerate}
\vspace{2mm}
\end{document}
